\section{Conclusion}

A genetic algorithm on NK fitness landscapes needs a substantially larger
population to reliably reach the global optimum as epistasis increases --
a real, sharp, population-size threshold. But the population-genetics
mechanism proposed to explain it, a single-locus selection coefficient
shrinking under epistasis until drift dominates (Kimura's classic
fixation-probability crossover), is falsified by direct measurement: the
selection coefficient measured from each landscape's distribution of
fitness effects grows, not shrinks, with epistasis, so the textbook
formula predicts the opposite trend from what is observed. Genetic
algorithms may be named for evolution, but at least this piece of their
population-sizing behaviour is not explained by evolution's own
single-locus theory of drift and selection -- a multi-peak exploration
argument is proposed as the more likely mechanism, left for future work to
test directly.
